\abstract{РЕФЕРАТ}

Объем работы равен \formbytotal{lastpage}{страниц}{е}{ам}{ам}. Работа содержит \formbytotal{figurecnt}{иллюстраци}{ю}{и}{й}, \formbytotal{tablecnt}{таблиц}{у}{ы}{} и \arabic{bibcount} библиографических источников. Количество приложений~--~1.

Перечень ключевых слов: интернет-провайдер, MikroTik, RouterOS, биллинг, автоматизация, веб-приложение, Laravel, Svelte, REST API, брандмауэр, простые очереди, разграничение доступа, тарифный план, клиент-серверная архитектура, реляционная база данных.

Объектом разработки является веб-система для автоматизированного управления клиентами интернет-провайдера малого и среднего масштаба, интегрированная с маршрутизаторами MikroTik по протоколу RouterOS~API. В состав системы входят серверная часть (программный интерфейс REST API) и панель администрирования для сотрудников оператора.

Цель работы~--~снижение трудозатрат сотрудников интернет-провайдера на рутинные операции по обслуживанию абонентов за счет автоматизации ключевых бизнес-процессов: подключения абонентов, формирования счетов, приостановки и возобновления услуг по факту оплаты.

В процессе работы проведен сравнительный анализ существующих систем биллинга, сформированы требования, спроектированы клиент-серверная архитектура и реляционная база данных, реализованы программный интерфейс в стиле REST, подсистема интеграции с MikroTik и подсистема автоматического биллинга с приостановкой и восстановлением обслуживания абонентов-должников.

При разработке использованы язык PHP с фреймворком Laravel, СУБД MySQL, библиотека для работы с RouterOS API, библиотека Svelte и метафреймворк SvelteKit.

Разработанная система прошла тестирование в стендовых условиях с подключенным маршрутизатором MikroTik и готова к опытной эксплуатации в небольших интернет-провайдерах.

\selectlanguage{english}
\abstract{ABSTRACT}

The volume of work is \formbytotal{lastpage}{page}{}{s}{s}. The work contains \formbytotal{figurecnt}{illustration}{}{s}{s}, \formbytotal{tablecnt}{table}{}{s}{s} and \arabic{bibcount} bibliographic sources. The number of appendices is~1.

List of keywords: Internet service provider, MikroTik, RouterOS, billing, automation, web application, Laravel, Svelte, REST API, firewall, simple queues, access control, service plan, client-server architecture, relational database.

The object of development is a web system for automated management of subscribers of a small and medium-scale Internet service provider, integrated with MikroTik routers via the RouterOS~API protocol. The system consists of a server-side component (REST API) and an administration panel for the operator's staff.

The purpose of the work is to reduce the labor costs of the provider's staff on routine subscriber service operations by automating the key business processes: subscriber connection, invoice generation, suspension and reactivation of services based on payment status, and network state monitoring.

During the work, a comparative analysis of existing billing and ISP management systems was performed, functional and non-functional requirements were formed, a multi-layer client-server architecture and a relational database were designed, and a REST-style API, a MikroTik integration subsystem and an automatic billing subsystem with suspension and reactivation of debtor subscribers were implemented.

The development used the PHP programming language within the Laravel framework, the MySQL relational DBMS, a library for the RouterOS API, and the Svelte library with the SvelteKit metaframework for building the client-side as a single-page application.

The developed system has been tested in bench conditions with a connected MikroTik router and is ready for pilot operation in small Internet service providers.

\selectlanguage{russian}
